% Test of code side by side (didactic) on the slides: the same pairs as
% test-pair.tex, in beamer frames.
%
% Build: see the Makefile.
%
% Expected in the resulting PDF:
%   1  Frame 1: the two main functions side by side, filling the frame, in a
%      small font (fontsize=\tiny), with the file names above.
%   2  Frame 2: \runbyrun, "Hello, world!" beside "hej", in the document's
%      size for minted, not the small size of code pairs.
%   3  Frame 3, not fragile: the same pair as frame 1.
%   4  Frame 4: the whole programs in the frame's normal size don't fit: the
%      lines run into the next column and off the frame, and the log has
%      one warning "A line of code side by side is ...pt wider than its
%      column ..., near page 4" (and one for frame 9); in a fragile frame the line
%      counts from the start of the frame.
%   5  Frame 5: the same pair with breaklines: the long lines wrap, no
%      warning.
%   6  Frame 6: the didacticpair of runs from the documentation, in the
%      size the frame sets with \setminted (footnotesize), not \tiny;
%      no warning.
%   7  Frame 7: slidewidth=frame, the main functions at \scriptsize: the
%      pair starts and ends half a gap (0.5em) from the slide's edges,
%      outside the text margins (WIDEL at about x=5pt, not 28pt).
%   8  Frame 8: the same pair in an example block: as wide as in frame 7,
%      beyond the block's background.
%   9  Frame 9: the same pair in the left one of two columns: within the
%      column, as without the option; it's too wide for the column at
%      \scriptsize, so the log warns "... near page 9".
%   The labels of frames 1, 4 and 5 are in the size of the code under them.
\documentclass{beamer}
\usepackage{didactic}
\begin{document}
\begin{frame}[fragile]
  \frametitle{Two versions}
  \codebycode[name=main,label]{python}{pairgreet1.py}%
    [name=main,label]{python}{pairgreet2.py}
\end{frame}
\begin{frame}[fragile]
  \frametitle{Two runs}
  \runbyrun{hello.py}[shell]{echo hej}
\end{frame}
\begin{frame}
  \frametitle{Not fragile}
  \codebycode[name=main]{python}{pairgreet1.py}[name=main]{python}{pairgreet2.py}
\end{frame}
\codebycodesetup{slides={fontsize=\normalsize}}
\begin{frame}[fragile]
  \frametitle{Too wide}
  \codebycode[label]{python}{pairgreet1.py}[label]{python}{pairgreet2.py}
\end{frame}
\begin{frame}[fragile]
  \frametitle{Too wide, wrapped}
  \codebycode[label,breaklines]{python}{pairgreet1.py}%
    [label,breaklines]{python}{pairgreet2.py}
\end{frame}
\begin{frame}[fragile]
  \frametitle{Runs by hand}
  \setminted{fontsize=\footnotesize}
  \begin{didacticpair}[runs]{2}
  \didacticpairlabel{\texttt{"w"}} & \didacticpairlabel{\texttt{"a"}} \\
  \runpython[transcript, args={w one}]{pairwrite.py}
  \runpython[transcript, args={w two}]{pairwrite.py}
  \runpython[transcript]{pairread.py}
  &
  \runpython[transcript, args={w one}]{pairwrite.py}
  \runpython[transcript, args={a two}]{pairwrite.py}
  \runpython[transcript]{pairread.py}
\end{didacticpair}

\end{frame}
\codebycodesetup{slidewidth=frame, slides={fontsize=\scriptsize}}
\begin{frame}[fragile]
  \frametitle{Slide width}
  \codebycode[name=main,label={WIDEL}]{python}{pairgreet1.py}%
    [name=main,label={WIDER}]{python}{pairgreet2.py}
\end{frame}
\begin{frame}[fragile]
  \begin{example}[In a block]
  \codebycode[name=main,label={BLOCKL}]{python}{pairgreet1.py}%
    [name=main,label={BLOCKR}]{python}{pairgreet2.py}
  \end{example}
\end{frame}
\begin{frame}[fragile]
  \begin{columns}
  \begin{column}{0.6\textwidth}
  \codebycode[name=main,label={COLL}]{python}{pairgreet1.py}%
    [name=main,label={COLR}]{python}{pairgreet2.py}
  \end{column}
  \begin{column}{0.35\textwidth}
  Text beside.
  \end{column}
  \end{columns}
\end{frame}
\end{document}
